\documentclass[aspectratio=169,10pt]{beamer}

\usetheme{SimplePlus}
\usepackage[T1]{fontenc}
\usepackage[utf8]{inputenc}
\usepackage{amsmath,amssymb}
\usepackage{graphicx}
\usepackage{booktabs}
\usepackage{hyperref}
\usepackage{enumitem}

\setlist[enumerate]{label=(\alph*)}

\title{Macroeconomics PhD Intro\\Practice Problems 1}
\subtitle{Solow Model \& Dynamic Optimization}
\author{Swapnil Singh}
\institute{Lietuvos Bankas \and Kaunas University of Technology}
\date{Spring 2026}

\newcommand{\apxbutton}[1]{\hfill\hyperlink{#1}{\beamerbutton{Solution}}}
\newcommand{\backbutton}[1]{\hfill\hyperlink{#1}{\beamerbutton{Back}}}

\begin{document}

\begin{frame}
  \titlepage
\end{frame}

\begin{frame}{Overview}
\begin{itemize}
  \item Problems cover material from Lectures 2 (Solow Model) and 3 (Dynamic Optimization).
  \item Organized in three tiers: \textbf{Easy}, \textbf{Hard}, and \textbf{Very Hard}.
  \item Mix of intuitive reasoning and analytical derivation.
  \item Each problem has a detailed solution accessible via the \beamerbutton{Solution} button.
\end{itemize}
\end{frame}

%% ============================================================
%% EASY PROBLEMS
%% ============================================================

\section{Easy Problems}

%% ----- EASY 1 -----
\begin{frame}{Easy 1: Steady State Interpretation}
\hypertarget{q-e1}{}
Consider the basic Solow model with Cobb-Douglas production $f(k)=k^\alpha$, saving rate $s$, and depreciation rate $\delta$.

\begin{enumerate}
  \item Write down the steady-state condition and solve for $\bar{k}$.
  \item Compute steady-state output $\bar{y}$, consumption $\bar{c}$, and investment $\bar{i}$.
  \item Using $\alpha=1/3$, $s=0.2$, $\delta=0.05$, compute the numerical values of $\bar{k}$, $\bar{y}$, and $\bar{k}/\bar{y}$.
  \item Explain in words: why does a higher $s$ raise the \emph{level} of $\bar{y}$ but not the long-run \emph{growth rate}?
\end{enumerate}
\apxbutton{s-e1}
\end{frame}

%% ----- EASY 2 -----
\begin{frame}{Easy 2: Capital-Output Ratio}
\hypertarget{q-e2}{}
In the Solow model, the steady-state capital-output ratio satisfies $\bar{k}/\bar{y} = s/\delta$.

\begin{enumerate}
  \item Show that this result holds for \emph{any} neoclassical production function $f(k)$, not just Cobb-Douglas.
  \item If $s = 0.25$ and $\delta = 0.05$, what is $\bar{k}/\bar{y}$? Is this consistent with U.S.\ data?
  \item Suppose depreciation doubles to $\delta = 0.10$ while $s$ stays fixed. What happens to $\bar{k}/\bar{y}$? Provide economic intuition.
\end{enumerate}
\apxbutton{s-e2}
\end{frame}

%% ----- EASY 3 -----
\begin{frame}{Easy 3: Convergence Intuition}
\hypertarget{q-e3}{}
Two countries, A and B, have identical parameters $(\alpha, s, \delta, \gamma, n)$ but different initial capital stocks: $k_0^A < k_0^B$.

\begin{enumerate}
  \item Which country grows faster in per-capita income during the transition? Why?
  \item Will they eventually reach the same level of per-capita income? The same growth rate?
  \item Does the Solow model predict unconditional or conditional convergence? What is the difference?
  \item How does this prediction fare against the data for all countries? For OECD countries?
\end{enumerate}
\apxbutton{s-e3}
\end{frame}

%% ----- EASY 4 -----
\begin{frame}{Easy 4: Euler Equation Intuition}
\hypertarget{q-e4}{}
The consumption Euler equation in the neoclassical growth model is:
\[
u'(c_t) = \beta\left[f'(k_{t+1}) + 1 - \delta\right] u'(c_{t+1})
\]

\begin{enumerate}
  \item Explain in words what this equation says about the household's intertemporal trade-off.
  \item What does $f'(k_{t+1}) + 1 - \delta$ represent economically?
  \item If $\beta(1+r) > 1$, is consumption growing or falling over time? Explain.
  \item At the steady state, what does the Euler equation imply about $f'(\bar{k})$?
\end{enumerate}
\apxbutton{s-e4}
\end{frame}

%% ----- EASY 5 -----
\begin{frame}{Easy 5: Sequential vs.\ Recursive}
\hypertarget{q-e5}{}
Consider the infinite-horizon neoclassical growth model.

\begin{enumerate}
  \item What is the ``state variable'' and what is the ``control variable'' in the planner's problem?
  \item Write down the Bellman equation for this problem.
  \item Why does infinite horizon require a transversality condition (TVC) but finite horizon does not?
  \item In the log-Cobb-Douglas case with $\delta=1$, the optimal policy is $k'=\alpha\beta Ak^\alpha$. What is the implied saving rate? How does this compare to the Solow model?
\end{enumerate}
\apxbutton{s-e5}
\end{frame}

%% ----- EASY 6 -----
\begin{frame}{Easy 6: CRRA and Balanced Growth}
\hypertarget{q-e6}{}
The CRRA utility function is $u(c) = \frac{c^{1-\sigma}-1}{1-\sigma}$.

\begin{enumerate}
  \item What happens as $\sigma \to 1$? Verify using L'H\^{o}pital's rule.
  \item What is the elasticity of intertemporal substitution (EIS) under CRRA?
  \item If $\sigma = 2$ and the interest rate rises by 1\%, by how much does consumption growth change?
  \item Why is CRRA the ``right'' utility class for models with balanced growth?
\end{enumerate}
\apxbutton{s-e6}
\end{frame}

%% ============================================================
%% HARD PROBLEMS
%% ============================================================

\section{Hard Problems}

%% ----- HARD 1 -----
\begin{frame}{Hard 1: Convergence Speed Derivation}
\hypertarget{q-h1}{}
In the Solow model with Cobb-Douglas production, no growth ($\gamma=n=0$):

\begin{enumerate}
  \item Linearize the fundamental equation $k_{t+1} = (1-\delta)k_t + sk_t^\alpha$ around the steady state $\bar{k}$.
  \item Show that the convergence speed is $\lambda = \delta(1-\alpha)$.
  \item Compute the half-life of convergence for $\alpha=1/3$, $\delta=0.05$.
  \item Explain why $s$ does not appear in $\lambda$ under Cobb-Douglas. What are the two opposing forces?
  \item For a general CES production function $f(k) = [ak^\rho + (1-a)]^{1/\rho}$, would $s$ affect $\lambda$? Explain.
\end{enumerate}
\apxbutton{s-h1}
\end{frame}

%% ----- HARD 2 -----
\begin{frame}{Hard 2: Growing Economy Steady State}
\hypertarget{q-h2}{}
Consider the Solow model with technology growth $\gamma$ and population growth $n$, Cobb-Douglas production.

\begin{enumerate}
  \item Derive the fundamental equation in efficiency units: $(1+\gamma)(1+n)\tilde{k}_{t+1} = (1-\delta)\tilde{k}_t + s\tilde{k}_t^\alpha$.
  \item Solve for the balanced growth path (BGP) value $\overline{\tilde{k}}$.
  \item Show that along the BGP, per-capita income grows at rate $\gamma$.
  \item Derive the convergence speed in the growing economy and show it equals $(1-\alpha)\left[1 - \frac{1-\delta}{(1+\gamma)(1+n)}\right]$.
  \item Using $\alpha=1/3$, $\gamma=0.02$, $n=0.01$, $\delta=0.046$, compute $\lambda$ and the half-life. Compare with the empirical estimate $\lambda \approx 0.02$.
\end{enumerate}
\apxbutton{s-h2}
\end{frame}

%% ----- HARD 3 -----
\begin{frame}{Hard 3: Log-Linearization of the Solow-RBC Model}
\hypertarget{q-h3}{}
Consider $k_{t+1} = sA_tk_t^\alpha\ell^{1-\alpha} + (1-\delta)k_t$ with persistent technology shocks $\hat{A}_t = \rho^t\varepsilon$.

\begin{enumerate}
  \item Define log-deviations $\hat{k}_t \equiv \log(k_t/\bar{k})$ and derive the log-linearized capital law:
  \[
  \hat{k}_{t+1} = (1 - \delta(1-\alpha))\hat{k}_t + \delta\hat{A}_t
  \]
  \item Solve this linear difference equation forward to obtain a closed form for $\hat{k}_{t+1}$.
  \item Derive the output deviation $\hat{y}_t = \hat{A}_t + \alpha\hat{k}_t$. Explain the two channels through which a technology shock affects output.
  \item Using $\delta=0.046$, $\alpha=1/3$, $\rho=0.9$, $\varepsilon=0.01$: compute $\hat{k}_1$, $\hat{k}_5$, and $\hat{y}_0$.
\end{enumerate}
\apxbutton{s-h3}
\end{frame}

%% ----- HARD 4 -----
\begin{frame}{Hard 4: Two-Period Consumption-Saving with Borrowing Constraint}
\hypertarget{q-h4}{}
A household lives two periods with $u(c)=\log c$, $\beta=0.95$, $r=0.05$, $a_0=0$, and no-borrowing constraint $a_1 \geq 0$.

\begin{enumerate}
  \item Suppose $w_0 = 10$, $w_1 = 10$. Write down the Lagrangian with the borrowing constraint. Solve for optimal $(c_0, c_1, a_1)$. Does the constraint bind?
  \item Now suppose $w_0 = 5$, $w_1 = 20$. Solve again. Does the borrowing constraint bind now? Explain the economics.
  \item For the constrained case, show that the Euler equation holds as an \emph{inequality}: $u'(c_0) > \beta(1+r)u'(c_1)$.
  \item What is the welfare cost of the borrowing constraint in case (b)? Compute the consumption-equivalent variation.
\end{enumerate}
\apxbutton{s-h4}
\end{frame}

%% ----- HARD 5 -----
\begin{frame}{Hard 5: Factor Shares and Euler's Theorem}
\hypertarget{q-h5}{}
Consider a CES production function $Y = [aK^\rho + (1-a)(AL)^\rho]^{1/\rho}$ with competitive factor pricing.

\begin{enumerate}
  \item Derive the capital share $\alpha_K \equiv r_t K_t / Y_t$ as a function of $\tilde{k} \equiv K/(AL)$ and $\rho$.
  \item Show that when $\rho = 0$ (Cobb-Douglas), the capital share is constant and equal to $a$.
  \item If $\rho < 0$ (elasticity of substitution $<1$), does capital deepening raise or lower the capital share? Explain.
  \item The U.S.\ labor share has declined from about 0.66 to 0.58 over 1980--2020. What does this imply about $\rho$ if the change is driven by capital deepening?
\end{enumerate}
\apxbutton{s-h5}
\end{frame}

%% ----- HARD 6 -----
\begin{frame}{Hard 6: Finite-Horizon NGM Closed Form}
\hypertarget{q-h6}{}
Consider the finite-horizon NGM with $u(c)=\log c$, $f(k)=Ak^\alpha$, $\delta=1$, horizon $T$.

\begin{enumerate}
  \item Write down the Euler equation and the terminal condition.
  \item Using backward induction from $k_{T+1}=0$, show that the optimal saving rate at date $t$ is:
  \[
  s_t \equiv \frac{k_{t+1}}{Ak_t^\alpha} = \alpha\beta\frac{1-(\alpha\beta)^{T-t}}{1-(\alpha\beta)^{T-t+1}}
  \]
  \item Show that as $T \to \infty$, $s_t \to \alpha\beta$ (the infinite-horizon result).
  \item Compute $s_T$ (saving rate in the last period). Is this intuitive?
  \item Plot or sketch $s_t$ as a function of $T-t$ (time remaining). What is the economic logic?
\end{enumerate}
\apxbutton{s-h6}
\end{frame}

%% ============================================================
%% VERY HARD PROBLEMS
%% ============================================================

\section{Very Hard Problems}

%% ----- VERY HARD 1 -----
\begin{frame}{Very Hard 1: Golden Rule and Dynamic Efficiency}
\hypertarget{q-v1}{}
In the Solow model with Cobb-Douglas production:

\begin{enumerate}
  \item The ``Golden Rule'' saving rate $s^{GR}$ maximizes steady-state consumption. Derive $s^{GR}$ and the corresponding $\bar{k}^{GR}$.
  \item Show that $s^{GR} = \alpha$. At the Golden Rule, what is the relationship between the net marginal product of capital and the depreciation rate?
  \item In the optimizing NGM (Cass-Koopmans), the steady-state satisfies $f'(\bar{k}) = 1/\beta - 1 + \delta$. Under what condition on $\beta$ does the optimizing economy choose the Golden Rule?
  \item An economy with $s > s^{GR}$ is called ``dynamically inefficient.'' Explain why everyone can be made better off by reducing $s$. Can the Cass-Koopmans economy ever be dynamically inefficient?
\end{enumerate}
\apxbutton{s-v1}
\end{frame}

%% ----- VERY HARD 2 -----
\begin{frame}{Very Hard 2: Poverty Trap and Big Push}
\hypertarget{q-v2}{}
Suppose the production function is $f(k) = Bk$ for $k \leq k^*$ and $f(k) = Bk^* + A(k-k^*)^\alpha$ for $k > k^*$, with $A > B > 0$ and $\alpha \in (0,1)$.

\begin{enumerate}
  \item Sketch the map $G(k) = (1-\delta)k + sf(k)$ against the 45-degree line. How many steady states can exist?
  \item Find conditions on $s$, $\delta$, $B$ such that a poverty-trap steady state exists below $k^*$.
  \item Suppose the economy is stuck at the poverty trap. What is the minimum ``big push'' (increase in $k$) required to escape?
  \item A policy temporarily raises $s$ to $\hat{s}$ for $\tau$ periods. Derive the condition on $(\hat{s}, \tau)$ that allows escape from the poverty trap.
\end{enumerate}
\apxbutton{s-v2}
\end{frame}

%% ----- VERY HARD 3 -----
\begin{frame}{Very Hard 3: Value Function Iteration}
\hypertarget{q-v3}{}
Consider the NGM with $u(c)=\log c$, $f(k)=Ak^\alpha$, $\delta=1$.

\begin{enumerate}
  \item Start with guess $V_0(k) = 0$. Compute $V_1(k)$ by solving the Bellman equation with $V_0$ as continuation value. Find the associated policy $g_1(k)$.
  \item Compute $V_2(k)$ using $V_1$ as continuation value. Find $g_2(k)$.
  \item Conjecture that $V_n(k) = E_n + F_n \log k$ for constants $E_n, F_n$. Verify this conjecture and derive the recursions for $E_n$ and $F_n$.
  \item Show that $F_n \to \frac{\alpha}{1-\alpha\beta}$ as $n \to \infty$, and that the limiting policy is $k' = \alpha\beta Ak^\alpha$.
\end{enumerate}
\apxbutton{s-v3}
\end{frame}

%% ----- VERY HARD 4 -----
\begin{frame}{Very Hard 4: Solow with Human Capital}
\hypertarget{q-v4}{}
Extend the Solow model: $Y_t = K_t^\alpha H_t^\phi (A_tL_t)^{1-\alpha-\phi}$ with $\alpha+\phi<1$. Physical capital accumulates as before with saving rate $s_K$. Human capital: $H_{t+1} = (1-\delta_H)H_t + s_H Y_t$.

\begin{enumerate}
  \item Define $\tilde{k}_t \equiv K_t/(A_tL_t)$ and $\tilde{h}_t \equiv H_t/(A_tL_t)$. Derive the two-dimensional system for $(\tilde{k}_{t+1}, \tilde{h}_{t+1})$.
  \item Solve for the BGP values $(\overline{\tilde{k}}, \overline{\tilde{h}})$ in terms of $s_K, s_H, \delta, \delta_H, \alpha, \phi$.
  \item Show that the convergence speed in this model is $\lambda = (1-\alpha-\phi)\delta_{\text{eff}}$, where $\delta_{\text{eff}}$ is the effective depreciation rate. Explain why human capital ($\phi > 0$) slows convergence.
  \item With $\alpha=1/3$ and $\phi=1/3$, compute $\lambda$ and the half-life. Compare with the empirical estimate $\lambda \approx 0.02$ and explain why this extended model better matches the data.
\end{enumerate}
\apxbutton{s-v4}
\end{frame}

%% ----- VERY HARD 5 -----
\begin{frame}{Very Hard 5: TVC, nPg, and Ponzi Schemes}
\hypertarget{q-v5}{}
Consider an infinite-horizon consumption-saving problem with constant $r$ and initial assets $a_0>0$.

\begin{enumerate}
  \item Show that without the no-Ponzi condition, the household can achieve unbounded utility by choosing $c_t = M$ for any large $M$, financed by ever-growing debt.
  \item With nPg imposed, derive the lifetime budget constraint: $\sum_{t=0}^\infty \frac{c_t}{(1+r)^t} \leq a_0(1+r)$.
  \item The TVC says $\lim_{t\to\infty} \beta^t u'(c_t) a_{t+1} = 0$. Show that if this fails---i.e., $\lim_{t\to\infty} \beta^t u'(c_t) a_{t+1} > 0$---the household is over-saving and can be made better off.
  \item With CRRA utility $u(c)=c^{1-\sigma}/(1-\sigma)$ and $\beta(1+r)=1$, show that TVC holds if and only if $a_0 > 0$ is finite.
\end{enumerate}
\apxbutton{s-v5}
\end{frame}

%% ----- VERY HARD 6 -----
\begin{frame}{Very Hard 6: Cobweb Dynamics Beyond Solow}
\hypertarget{q-v6}{}
Consider the one-dimensional dynamical system $x_{t+1} = G(x_t)$ with a unique interior fixed point $x^*$ satisfying $G(x^*)=x^*$.

\begin{enumerate}
  \item Show that if $G'(x^*) \in (0,1)$, convergence is monotone. If $G'(x^*) \in (-1,0)$, convergence is oscillatory.
  \item For the Solow model $G(k)=(1-\delta)k+sk^\alpha$, show that $G'(\bar{k})=1-\delta(1-\alpha) \in (0,1)$ for all standard parameter values. Conclude that oscillatory convergence is impossible in the basic Solow model.
  \item Now consider a ``bakery production function'' where output declines at very high $k$: $f(k)=k^\alpha e^{-\eta k}$ for $\eta>0$ small. Find the condition on parameters such that $G'(\bar{k})<0$, generating oscillatory dynamics.
  \item For the bakery case, find the condition for chaos ($|G'(\bar{k})|>1$). Is this empirically plausible?
\end{enumerate}
\apxbutton{s-v6}
\end{frame}

%% ============================================================
%%                    SOLUTIONS
%% ============================================================

\appendix
\section{Solutions}

%% ----- SOLUTION EASY 1 -----
\begin{frame}{Solution: Easy 1 (a)--(b)}
\hypertarget{s-e1}{}

\textbf{(a)} Steady state: $\bar{k} = (1-\delta)\bar{k} + s\bar{k}^\alpha \implies \delta\bar{k} = s\bar{k}^\alpha$.
\[
\bar{k}^{1-\alpha} = \frac{s}{\delta} \implies \boxed{\bar{k} = \left(\frac{s}{\delta}\right)^{\frac{1}{1-\alpha}}}
\]

\textbf{(b)} $\bar{y} = \bar{k}^\alpha = \left(\frac{s}{\delta}\right)^{\frac{\alpha}{1-\alpha}}$. Investment: $\bar{i} = s\bar{y} = \delta\bar{k}$. Consumption:
\[
\bar{c} = (1-s)\bar{y} = (1-s)\left(\frac{s}{\delta}\right)^{\frac{\alpha}{1-\alpha}}
\]
\end{frame}

\begin{frame}{Solution: Easy 1 (c)--(d)}
\textbf{(c)} With $\alpha=1/3$, $s=0.2$, $\delta=0.05$:
\[
\bar{k} = (0.2/0.05)^{1/(1-1/3)} = 4^{3/2} = 8, \quad \bar{y} = 8^{1/3} = 2, \quad \bar{k}/\bar{y} = 4
\]

\textbf{(d)} A higher $s$ shifts the investment curve $sf(k)$ upward, leading to a higher intersection with $\delta k$---thus a higher $\bar{k}$ and $\bar{y}$. But in the long run, $\tilde{k}$ stabilizes and per-capita growth equals $\gamma$ (the exogenous technology growth rate). Saving determines the \textbf{level} of the balanced growth path, not its slope. Diminishing returns to capital ensure that extra investment cannot sustain permanently higher growth.
\backbutton{q-e1}
\end{frame}

%% ----- SOLUTION EASY 2 -----
\begin{frame}{Solution: Easy 2}
\hypertarget{s-e2}{}

\textbf{(a)} The steady-state condition for \emph{any} $f$ is $\delta\bar{k} = sf(\bar{k}) = s\bar{y}$. Rearranging: $\bar{k}/\bar{y} = s/\delta$. No functional-form assumption is used.

\medskip
\textbf{(b)} $\bar{k}/\bar{y} = 0.25/0.05 = 5$. U.S.\ data suggests $K/Y \approx 3$--$4$, so this is slightly high but broadly consistent.

\medskip
\textbf{(c)} If $\delta$ doubles: $\bar{k}/\bar{y} = 0.25/0.10 = 2.5$. The ratio halves. \textbf{Intuition:} faster depreciation erodes the capital stock more quickly. To maintain the same $K/Y$ ratio with higher $\delta$, you would need proportionally more saving. With fixed $s$, the ratio falls.
\backbutton{q-e2}
\end{frame}

%% ----- SOLUTION EASY 3 -----
\begin{frame}{Solution: Easy 3}
\hypertarget{s-e3}{}

\textbf{(a)} Country A (lower $k_0$) grows faster. With identical parameters, both share the same $\bar{k}$. Country A is further below $\bar{k}$, so diminishing returns imply a higher marginal product of capital and hence faster growth.

\medskip
\textbf{(b)} Yes to both. Identical parameters $\implies$ identical BGP. They converge to the same $\bar{k}$ (level) and the same growth rate $\gamma$ (growth).

\medskip
\textbf{(c)} \textbf{Conditional convergence}: countries with the same structural parameters converge to the same BGP. \textbf{Unconditional convergence}: all countries converge regardless of parameters. Solow predicts conditional, not unconditional.

\medskip
\textbf{(d)} All countries: no systematic relationship between initial income and growth (no unconditional convergence). OECD countries: clear negative relationship---conditional convergence holds among similar economies.
\backbutton{q-e3}
\end{frame}

%% ----- SOLUTION EASY 4 -----
\begin{frame}{Solution: Easy 4}
\hypertarget{s-e4}{}

\textbf{(a)} The marginal utility cost of giving up one unit of consumption today ($u'(c_t)$) must equal the discounted marginal utility gain from investing that unit, earning return $f'(k_{t+1})+1-\delta$, and consuming the proceeds tomorrow ($\beta[f'(k_{t+1})+1-\delta]u'(c_{t+1})$). At optimum, the household is indifferent between consuming now and later.

\medskip
\textbf{(b)} It is the gross return on capital: one unit invested yields $f'(k_{t+1})$ in marginal product plus $1-\delta$ in surviving capital.

\medskip
\textbf{(c)} If $\beta(1+r)>1$: the return to waiting exceeds impatience, so $u'(c_{t+1})<u'(c_t)$, meaning $c_{t+1}>c_t$---consumption is \textbf{growing}.

\medskip
\textbf{(d)} At steady state $c_{t+1}=c_t$, so $u'$ cancels: $1 = \beta[f'(\bar{k})+1-\delta]$, giving $f'(\bar{k})=1/\beta - 1+\delta$.
\backbutton{q-e4}
\end{frame}

%% ----- SOLUTION EASY 5 -----
\begin{frame}{Solution: Easy 5}
\hypertarget{s-e5}{}

\textbf{(a)} State: $k_t$ (capital, predetermined). Control: $c_t$ or equivalently $k_{t+1}$ (consumption/investment choice).

\medskip
\textbf{(b)} $V(k) = \max_{k' \in [0, f(k)+(1-\delta)k]} \left\{u(f(k)+(1-\delta)k - k') + \beta V(k')\right\}$.

\medskip
\textbf{(c)} Finite horizon has a terminal condition ($k_{T+1}=0$) that pins down the boundary. Infinite horizon has no terminal date, so TVC replaces it---ensuring the household doesn't accumulate assets with positive shadow value forever.

\medskip
\textbf{(d)} Saving rate: $s = k'/(Ak^\alpha) = \alpha\beta$. In the Solow model, $s$ is exogenous. Here it is endogenously pinned down by preferences ($\beta$) and technology ($\alpha$). With $\alpha=1/3$ and $\beta=0.96$, $s \approx 0.32$.
\backbutton{q-e5}
\end{frame}

%% ----- SOLUTION EASY 6 -----
\begin{frame}{Solution: Easy 6}
\hypertarget{s-e6}{}

\textbf{(a)} Write $u(c)=(c^{1-\sigma}-1)/(1-\sigma)$. As $\sigma\to 1$: apply L'H\^{o}pital by differentiating numerator and denominator with respect to $\sigma$:
\[
\lim_{\sigma\to 1}\frac{c^{1-\sigma}-1}{1-\sigma} = \lim_{\sigma\to 1}\frac{-c^{1-\sigma}\ln c}{-1} = \ln c
\]

\textbf{(b)} $\text{EIS} = 1/\sigma$. It measures the percent change in the consumption ratio $c_{t+1}/c_t$ per percent change in the gross return.

\medskip
\textbf{(c)} With $\sigma=2$: EIS $= 0.5$. A 1\% rise in the interest rate raises consumption growth by $0.5\%$.

\medskip
\textbf{(d)} Balanced growth requires that as $c_t$ grows at rate $\gamma$, the saving rate remains constant. CRRA is the unique class where the income and substitution effects of growing consumption on saving exactly offset, keeping $s$ constant along a balanced growth path.
\backbutton{q-e6}
\end{frame}

%% ----- SOLUTION HARD 1 -----
\begin{frame}{Solution: Hard 1 (a)--(b)}
\hypertarget{s-h1}{}

\textbf{(a)} Define $G(k) = (1-\delta)k + sk^\alpha$. Taylor expand around $\bar{k}$:
\[
k_{t+1} - \bar{k} \approx G'(\bar{k})(k_t - \bar{k})
\]
\[
G'(\bar{k}) = 1 - \delta + s\alpha\bar{k}^{\alpha-1}
\]

\textbf{(b)} From steady state: $s\bar{k}^\alpha = \delta\bar{k}$, so $s\bar{k}^{\alpha-1} = \delta$. Therefore $s\alpha\bar{k}^{\alpha-1} = \alpha\delta$. Substituting:
\[
G'(\bar{k}) = 1 - \delta + \alpha\delta = 1 - \delta(1-\alpha)
\]
The convergence speed is $\lambda \equiv 1 - G'(\bar{k}) = \delta(1-\alpha)$.
\end{frame}

\begin{frame}{Solution: Hard 1 (c)--(e)}
\textbf{(c)} $\lambda = 0.05 \times (1-1/3) = 0.05 \times 2/3 \approx 0.033$. Half-life: $h = \ln 2 / \lambda \approx 0.693/0.033 \approx 21$ years.

\medskip
\textbf{(d)} Two opposing forces from raising $s$:
\begin{enumerate}
  \item Higher $s$ increases the slope of $G$ (steeper investment curve) $\implies$ slower convergence.
  \item Higher $s$ raises $\bar{k}$, where $f'(\bar{k})$ is smaller (diminishing returns) $\implies$ faster convergence.
\end{enumerate}
With Cobb-Douglas, these exactly cancel: $sf'(\bar{k}) = s\alpha\bar{k}^{\alpha-1} = \alpha\delta$, independent of $s$.

\medskip
\textbf{(e)} For CES with $\rho \neq 0$, the elasticity of substitution differs from 1, so the cancellation no longer holds. The marginal product $f'(\bar{k})$ depends on $s$ through a nonlinear channel, and $s$ generally affects $\lambda$.
\backbutton{q-h1}
\end{frame}

%% ----- SOLUTION HARD 2 -----
\begin{frame}{Solution: Hard 2 (a)--(b)}
\hypertarget{s-h2}{}

\textbf{(a)} Start from $K_{t+1} = (1-\delta)K_t + sK_t^\alpha(A_tL_t)^{1-\alpha}$. Divide by $A_tL_t$:
\[
\frac{K_{t+1}}{A_tL_t} = (1-\delta)\tilde{k}_t + s\tilde{k}_t^\alpha
\]
The LHS: $\frac{K_{t+1}}{A_tL_t} = \frac{A_{t+1}L_{t+1}}{A_tL_t}\cdot\frac{K_{t+1}}{A_{t+1}L_{t+1}} = (1+\gamma)(1+n)\tilde{k}_{t+1}$.

\medskip
\textbf{(b)} Set $\tilde{k}_{t+1}=\tilde{k}_t=\overline{\tilde{k}}$:
\[
[(1+\gamma)(1+n)-(1-\delta)]\overline{\tilde{k}} = s\overline{\tilde{k}}^\alpha
\]
\[
\overline{\tilde{k}} = \left(\frac{s}{(1+\gamma)(1+n)+\delta-1}\right)^{\frac{1}{1-\alpha}}
\]
\end{frame}

\begin{frame}{Solution: Hard 2 (c)--(e)}
\textbf{(c)} Per-capita income: $y_t = A_t f(\tilde{k}_t)$. On BGP, $f(\tilde{k}_t) = f(\overline{\tilde{k}})$ constant. So $y_{t+1}/y_t = A_{t+1}/A_t = 1+\gamma$.

\medskip
\textbf{(d)} Linearize around BGP. The map is $\tilde{k}_{t+1} = \frac{(1-\delta)\tilde{k}_t + s\tilde{k}_t^\alpha}{(1+\gamma)(1+n)}$. The derivative at BGP:
\[
G'(\overline{\tilde{k}}) = \frac{1-\delta+s\alpha\overline{\tilde{k}}^{\alpha-1}}{(1+\gamma)(1+n)} = \frac{1-\delta+\alpha[(1+\gamma)(1+n)-1+\delta]}{(1+\gamma)(1+n)}
\]
\[
\lambda = 1 - G'(\overline{\tilde{k}}) = (1-\alpha)\left[1-\frac{1-\delta}{(1+\gamma)(1+n)}\right]
\]

\textbf{(e)} $\lambda = (2/3)[1 - 0.954/(1.02\times1.01)] = (2/3)[1 - 0.954/1.0302] = (2/3)(1-0.926) \approx 0.049$. Half-life $\approx 14$ years. Empirical $\lambda\approx0.02$ gives half-life $\approx 35$ years. The model over-predicts convergence speed by a factor of $\sim 2.5$.
\backbutton{q-h2}
\end{frame}

%% ----- SOLUTION HARD 3 -----
\begin{frame}{Solution: Hard 3 (a)--(b)}
\hypertarget{s-h3}{}

\textbf{(a)} Write $k_t = \bar{k}e^{\hat{k}_t}$, $A_t = \bar{A}e^{\hat{A}_t}$. Substitute into $k_{t+1}=sA_tk_t^\alpha\ell^{1-\alpha}+(1-\delta)k_t$:
\[
\bar{k}(1+\hat{k}_{t+1}) \approx s\bar{A}\bar{k}^\alpha\ell^{1-\alpha}(1+\hat{A}_t+\alpha\hat{k}_t) + (1-\delta)\bar{k}(1+\hat{k}_t)
\]
Using $\delta\bar{k} = s\bar{A}\bar{k}^\alpha\ell^{1-\alpha}$ and collecting first-order terms:
\[
\hat{k}_{t+1} = (1-\delta(1-\alpha))\hat{k}_t + \delta\hat{A}_t
\]

\textbf{(b)} With $\hat{A}_t = \rho^t\varepsilon$, iterate:
\[
\hat{k}_{t+1} = \delta\varepsilon\sum_{\tau=0}^{t}(1-\lambda)^\tau \rho^{t-\tau} = \delta\varepsilon\rho^t\frac{1-\left(\frac{1-\lambda}{\rho}\right)^{t+1}}{1-\frac{1-\lambda}{\rho}}
\]
where $\lambda = \delta(1-\alpha)$.
\end{frame}

\begin{frame}{Solution: Hard 3 (c)--(d)}
\textbf{(c)} From $y_t = A_tk_t^\alpha\ell^{1-\alpha}$, log-linearize: $\hat{y}_t = \hat{A}_t + \alpha\hat{k}_t$.

Two channels:
\begin{itemize}
  \item \textbf{Direct effect} ($\hat{A}_t$): technology shock raises output immediately.
  \item \textbf{Propagation effect} ($\alpha\hat{k}_t$): higher output $\to$ higher investment $\to$ higher future capital $\to$ higher future output. This channel is slower and more persistent.
\end{itemize}

\textbf{(d)} With $\delta=0.046$, $\alpha=1/3$, $\lambda=0.046\times 2/3\approx 0.031$:
\begin{itemize}
  \item $\hat{k}_1 = \delta\varepsilon = 0.046\times 0.01 = 0.00046$ (0.046\%).
  \item $\hat{k}_5 = 0.046\times 0.01\times 0.9^4 \frac{1-(0.969/0.9)^5}{1-0.969/0.9} \approx 0.00198$ (0.198\%).
  \item $\hat{y}_0 = \hat{A}_0 + \alpha\hat{k}_0 = 0.01 + 0 = 0.01$ (1\%, since $\hat{k}_0=0$).
\end{itemize}
\backbutton{q-h3}
\end{frame}

%% ----- SOLUTION HARD 4 -----
\begin{frame}{Solution: Hard 4 (a)--(b)}
\hypertarget{s-h4}{}

\textbf{(a)} $w_0=w_1=10$. Lagrangian: $\mathcal{L} = \log c_0 + 0.95\log c_1 + \mu(10+10/1.05-c_0-c_1/1.05) + \lambda(10-c_0)$.

Try interior ($\lambda=0$): Euler gives $1/c_0 = 0.95\times 1.05/c_1$, so $c_1 = 0.9975\,c_0$. Budget: $c_0 + c_1/1.05 = 10+10/1.05 = 19.524$. Solving: $c_0 \approx 10.013$, $c_1 \approx 9.988$, $a_1 = 10-c_0 \approx -0.013$. Since $a_1 < 0$, constraint binds!

Re-solve with $a_1=0$: $c_0=10$, $c_1=10$. Constraint binds but barely.

\medskip
\textbf{(b)} $w_0=5$, $w_1=20$. Interior: $c_1=0.9975\,c_0$. Budget: $c_0+c_1/1.05 = 5+20/1.05=24.048$. Solving: $c_0\approx 12.33$. But $a_1 = 5-12.33<0$. Constraint binds hard: $c_0=5$, $c_1=20$.

Economics: the household wants to borrow against high future income but cannot.
\end{frame}

\begin{frame}{Solution: Hard 4 (c)--(d)}
\textbf{(c)} At the constrained optimum $c_0=5$, $c_1=20$:
\[
u'(c_0) = 1/5 = 0.2, \quad \beta(1+r)u'(c_1) = 0.95\times 1.05\times(1/20) = 0.04988
\]
Since $0.2 > 0.04988$: $u'(c_0) > \beta(1+r)u'(c_1)$. The household would like to borrow (shift consumption forward) but cannot---marginal utility of current consumption exceeds the discounted return from saving.

\medskip
\textbf{(d)} Unconstrained optimum: $c_0^*\approx 12.33$, $c_1^*\approx 12.30$. Constrained: $c_0^c=5$, $c_1^c=20$.

Welfare comparison: $U^* = \log 12.33 + 0.95\log 12.30 \approx 4.887$. $U^c = \log 5 + 0.95\log 20 \approx 4.457$.

Consumption-equivalent variation $\xi$: $(1+\xi)$ solves $\log(5(1+\xi))+0.95\log(20(1+\xi)) = U^*$. This gives $1.95\log(1+\xi) = 0.43$, so $\xi \approx e^{0.22}-1 \approx 24.6\%$.
\backbutton{q-h4}
\end{frame}

%% ----- SOLUTION HARD 5 -----
\begin{frame}{Solution: Hard 5 (a)--(b)}
\hypertarget{s-h5}{}

\textbf{(a)} $Y = [aK^\rho + (1-a)(AL)^\rho]^{1/\rho}$. Competitive pricing: $r = \partial Y/\partial K = a\left(\frac{Y}{K}\right)^{1-\rho}$. Capital share:
\[
\alpha_K = \frac{rK}{Y} = a\left(\frac{K}{Y}\right)^{\rho} \cdot \frac{Y}{K}\cdot K/Y = a\left(\frac{K}{Y}\right)^{\rho}
\]
More precisely:
\[
\alpha_K = \frac{a K^\rho}{aK^\rho + (1-a)(AL)^\rho} = \frac{a\tilde{k}^\rho}{a\tilde{k}^\rho + (1-a)}
\]

\textbf{(b)} As $\rho\to 0$, take logs: $\alpha_K = \frac{a\tilde{k}^0}{a\tilde{k}^0+(1-a)} = \frac{a}{a+(1-a)} = a$. Constant, independent of $\tilde{k}$.
\end{frame}

\begin{frame}{Solution: Hard 5 (c)--(d)}
\textbf{(c)} Differentiate $\alpha_K$ w.r.t.\ $\tilde{k}$:
\[
\frac{\partial\alpha_K}{\partial\tilde{k}} = \frac{a\rho\tilde{k}^{\rho-1}(1-a)}{[a\tilde{k}^\rho+(1-a)]^2}
\]
The sign is $\text{sign}(\rho)$. If $\rho<0$ (elasticity of substitution $\sigma_s = 1/(1-\rho)<1$): capital deepening ($\tilde{k}\uparrow$) \emph{lowers} the capital share, hence \emph{raises} the labor share.

Wait---but the labor share has \emph{declined}. So with $\rho<0$ and capital deepening, we'd predict \emph{rising} labor share. Contradiction.

\medskip
\textbf{(d)} The declining U.S.\ labor share with capital deepening requires $\rho>0$ (elasticity $>1$). Capital deepening raises the capital share when capital and labor are gross substitutes. Recent estimates (Karabarbounis and Neiman, 2014) indeed argue for $\sigma_s > 1$, driven by falling relative price of investment goods.
\backbutton{q-h5}
\end{frame}

%% ----- SOLUTION HARD 6 -----
\begin{frame}{Solution: Hard 6 (a)--(c)}
\hypertarget{s-h6}{}

\textbf{(a)} Euler: $\frac{1}{c_t} = \beta\frac{\alpha Ak_{t+1}^{\alpha-1}}{c_{t+1}}$. Terminal: $k_{T+1}=0$ (no value to leaving capital).

\medskip
\textbf{(b)} At $t=T$: $k_{T+1}=0$, so $c_T=Ak_T^\alpha$ and $s_T=0$. At $t=T-1$: use Euler to get $k_T = \alpha\beta Ak_{T-1}^\alpha/(1+\alpha\beta)$, giving $s_{T-1} = \alpha\beta/(1+\alpha\beta)$. By induction:
\[
s_t = \alpha\beta\frac{1-(\alpha\beta)^{T-t}}{1-(\alpha\beta)^{T-t+1}}
\]

\textbf{(c)} As $T-t \to \infty$: $(\alpha\beta)^{T-t}\to 0$ (since $\alpha\beta<1$), so $s_t \to \alpha\beta\frac{1}{1} = \alpha\beta$.
\end{frame}

\begin{frame}{Solution: Hard 6 (d)--(e)}
\textbf{(d)} $s_T = \alpha\beta\frac{1-(\alpha\beta)^0}{1-(\alpha\beta)^1} = \alpha\beta\frac{0}{1-\alpha\beta} = 0$.

This is completely intuitive: in the last period, there is no future to save for. All output is consumed.

\medskip
\textbf{(e)} As a function of ``time remaining'' $\tau \equiv T-t$:
\begin{itemize}
  \item $s(\tau=0) = 0$: no saving at the end.
  \item $s(\tau)$ increases monotonically toward $\alpha\beta$.
  \item Convergence is geometric at rate $\alpha\beta$.
\end{itemize}

\textbf{Economic logic:} With more periods ahead, the return to investing (measured by total future output produced) is higher, so the household saves more. The infinite-horizon limit represents the case where the investment motive is strongest, yielding the maximal constant saving rate $\alpha\beta$.
\backbutton{q-h6}
\end{frame}

%% ----- SOLUTION VERY HARD 1 -----
\begin{frame}{Solution: Very Hard 1 (a)--(b)}
\hypertarget{s-v1}{}

\textbf{(a)} Steady-state consumption: $\bar{c} = f(\bar{k})-\delta\bar{k}$. From the steady-state condition $\delta\bar{k}=s\bar{k}^\alpha$:
\[
\bar{c}(s) = \bar{k}^\alpha - \delta\bar{k} = \left(\frac{s}{\delta}\right)^{\frac{\alpha}{1-\alpha}} - \delta\left(\frac{s}{\delta}\right)^{\frac{1}{1-\alpha}}
\]
FOC: $\frac{d\bar{c}}{ds}=0$ yields $f'(\bar{k}^{GR})\frac{d\bar{k}}{ds} = \delta\frac{d\bar{k}}{ds}$, so $f'(\bar{k}^{GR})=\delta$.

With Cobb-Douglas: $\alpha\bar{k}^{\alpha-1}=\delta \implies \bar{k}^{GR} = (\alpha/\delta)^{1/(1-\alpha)}$.

\medskip
\textbf{(b)} Compare $\bar{k}^{GR} = (\alpha/\delta)^{1/(1-\alpha)}$ with $\bar{k}=(s/\delta)^{1/(1-\alpha)}$: Golden Rule requires $s^{GR}=\alpha$. At this point $f'(\bar{k}^{GR})=\delta$, meaning the net marginal product of capital equals zero (gross MPK equals the depreciation rate).
\end{frame}

\begin{frame}{Solution: Very Hard 1 (c)--(d)}
\textbf{(c)} In Cass-Koopmans: $f'(\bar{k})=1/\beta-1+\delta$. Golden Rule requires $f'(\bar{k}^{GR})=\delta$. These coincide when:
\[
1/\beta-1+\delta=\delta \implies \beta=1
\]
Only with no discounting ($\beta=1$) does the optimizing economy choose the Golden Rule. With $\beta<1$, impatience leads to $\bar{k}<\bar{k}^{GR}$---the economy under-accumulates relative to Golden Rule.

\medskip
\textbf{(d)} Dynamic inefficiency ($s>s^{GR}$) means $f'(\bar{k})<\delta$: the marginal product of capital is below depreciation. Reducing capital frees resources for consumption in every period (consuming the excess capital and reducing replacement investment). Everyone gains---a Pareto improvement.

The Cass-Koopmans economy \textbf{cannot} be dynamically inefficient because $\beta<1$ implies $f'(\bar{k})=1/\beta-1+\delta>\delta$. Optimizing agents never over-accumulate.
\backbutton{q-v1}
\end{frame}

%% ----- SOLUTION VERY HARD 2 -----
\begin{frame}{Solution: Very Hard 2 (a)--(b)}
\hypertarget{s-v2}{}

\textbf{(a)} For $k\leq k^*$: $G(k) = (1-\delta)k + sBk$. This is linear with slope $(1-\delta+sB)$. For $k>k^*$: $G$ transitions to a concave function (Solow-like with $A>\!B$). The map can cross the 45-degree line up to three times: once in the linear region (if $sB<\delta$), once near $k^*$ (unstable), and once in the concave region (stable, high steady state).

\medskip
\textbf{(b)} Poverty-trap steady state in the linear region requires $G(k)=k$ for some $k<k^*$:
\[
(1-\delta+sB)k = k \implies sB = \delta
\]
But this gives $G(k)=k$ for \emph{all} $k\leq k^*$. For a proper trap at $k<k^*$ (where $G$ crosses from above): we need $sB<\delta$. Then $G(k)<k$ in the linear region and the only low steady state is $\bar{k}_1=0$. For a non-trivial poverty trap, we need the production function to have a region of increasing returns near $k^*$ so that $G$ crosses the 45-degree line from below at some positive $\bar{k}_1>0$.

With $sB<\delta$: the poverty trap is at $\bar{k}_1=0$.
\end{frame}

\begin{frame}{Solution: Very Hard 2 (c)--(d)}
\textbf{(c)} The economy escapes the trap by reaching a capital level above the unstable steady state $\bar{k}_2$ (where $G$ crosses the 45-degree line from below to above). The minimum big push is $\Delta k = \bar{k}_2 - k_{\text{trap}}$.

For the piecewise specification: $\bar{k}_2$ is near $k^*$ where the transition from linear to concave production occurs. Precisely, $\bar{k}_2$ solves $G(\bar{k}_2)=\bar{k}_2$ in the increasing-returns region.

\medskip
\textbf{(d)} Starting from $k_0 = \bar{k}_1$ (the trap), temporarily raise $s$ to $\hat{s}$ for $\tau$ periods. Need $k_\tau > \bar{k}_2$ after $\tau$ periods of iteration under $\hat{s}$. In the linear region, $k_t = (1-\delta+\hat{s}B)^t k_0$ when $\hat{s}B>\delta$. The condition is:
\[
(1-\delta+\hat{s}B)^\tau \cdot k_0 > \bar{k}_2
\]
\[
\tau > \frac{\ln(\bar{k}_2/k_0)}{\ln(1-\delta+\hat{s}B)}
\]
Higher $\hat{s}$ allows escape in fewer periods.
\backbutton{q-v2}
\end{frame}

%% ----- SOLUTION VERY HARD 3 -----
\begin{frame}{Solution: Very Hard 3 (a)--(b)}
\hypertarget{s-v3}{}

\textbf{(a)} With $V_0(k)=0$: the Bellman becomes $V_1(k) = \max_{k'\geq 0}\log(Ak^\alpha - k')$. This is maximized at $k'=0$ (no continuation value to investing). So $g_1(k)=0$ and
\[
V_1(k) = \log(Ak^\alpha) = \log A + \alpha\log k
\]

\textbf{(b)} $V_2(k) = \max_{k'\geq 0}\{\log(Ak^\alpha - k') + \beta[\log A + \alpha\log k']\}$. FOC:
\[
\frac{1}{Ak^\alpha - k'} = \frac{\beta\alpha}{k'} \implies k' = \frac{\alpha\beta}{1+\alpha\beta}Ak^\alpha
\]
So $g_2(k) = \frac{\alpha\beta}{1+\alpha\beta}Ak^\alpha$ and:
\[
V_2(k) = (1+\alpha\beta)\log\frac{Ak^\alpha}{1+\alpha\beta} + \alpha\beta\log\frac{\alpha\beta Ak^\alpha}{1+\alpha\beta}
\]
which has the form $E_2 + F_2\log k$ with $F_2 = \alpha(1+\alpha\beta)$.
\end{frame}

\begin{frame}{Solution: Very Hard 3 (c)--(d)}
\textbf{(c)} Conjecture $V_n(k) = E_n + F_n\log k$. Bellman gives:
\[
V_{n+1}(k) = \max_{k'}\{\log(Ak^\alpha - k') + \beta(E_n + F_n\log k')\}
\]
FOC: $k' = \frac{\beta F_n}{1+\beta F_n}Ak^\alpha$. Substituting back yields:
\[
F_{n+1} = \alpha(1+\beta F_n)
\]
$E_{n+1}$ follows a similar recursion (involving $\log A$ and $F_n$ terms).

\medskip
\textbf{(d)} Steady state of $F_{n+1}=\alpha(1+\beta F_n)$: $F^* = \alpha + \alpha\beta F^*$, so:
\[
F^* = \frac{\alpha}{1-\alpha\beta}
\]
The policy function converges to $k' = \frac{\beta F^*}{1+\beta F^*}Ak^\alpha = \frac{\alpha\beta}{1-\alpha\beta+\alpha\beta}Ak^\alpha = \alpha\beta Ak^\alpha$.

This matches the sequential solution exactly, confirming equivalence.
\backbutton{q-v3}
\end{frame}

%% ----- SOLUTION VERY HARD 4 -----
\begin{frame}{Solution: Very Hard 4 (a)--(b)}
\hypertarget{s-v4}{}

\textbf{(a)} Divide $K_{t+1}=(1-\delta)K_t+s_KY_t$ and $H_{t+1}=(1-\delta_H)H_t+s_HY_t$ by $A_{t+1}L_{t+1} = (1+\gamma)(1+n)A_tL_t$. Output per efficiency unit: $\tilde{y}_t = \tilde{k}_t^\alpha\tilde{h}_t^\phi$. The system is:
\begin{align*}
(1+\gamma)(1+n)\tilde{k}_{t+1} &= (1-\delta)\tilde{k}_t + s_K\tilde{k}_t^\alpha\tilde{h}_t^\phi \\
(1+\gamma)(1+n)\tilde{h}_{t+1} &= (1-\delta_H)\tilde{h}_t + s_H\tilde{k}_t^\alpha\tilde{h}_t^\phi
\end{align*}

\textbf{(b)} At BGP, set $\tilde{k}_{t+1}=\tilde{k}_t$, $\tilde{h}_{t+1}=\tilde{h}_t$. Let $D\equiv(1+\gamma)(1+n)-1+\delta$ and $D_H\equiv(1+\gamma)(1+n)-1+\delta_H$:
\[
D\,\overline{\tilde{k}} = s_K\overline{\tilde{k}}^\alpha\overline{\tilde{h}}^\phi, \qquad D_H\,\overline{\tilde{h}} = s_H\overline{\tilde{k}}^\alpha\overline{\tilde{h}}^\phi
\]
Dividing: $\overline{\tilde{h}}/\overline{\tilde{k}} = (s_H D)/(s_K D_H)$. Substituting back yields closed-form expressions for both.
\end{frame}

\begin{frame}{Solution: Very Hard 4 (c)--(d)}
\textbf{(c)} Linearizing the two-dimensional system around BGP, the eigenvalues of the Jacobian govern convergence. The dominant eigenvalue (closest to 1) determines the speed. By analogy with the one-dimensional case, the convergence speed generalizes to:
\[
\lambda \approx (1-\alpha-\phi)\left[1-\frac{1-\delta_{\text{eff}}}{(1+\gamma)(1+n)}\right]
\]
where $\delta_{\text{eff}}$ is a weighted average of $\delta$ and $\delta_H$. The key is that $\alpha+\phi$ replaces $\alpha$: human capital acts like a second accumulable factor, strengthening diminishing-returns attenuation and slowing convergence.

\medskip
\textbf{(d)} With $\alpha=\phi=1/3$: $1-\alpha-\phi = 1/3$, roughly a third of the baseline Solow value. Using similar $\delta$ values:
\[
\lambda \approx (1/3)\times 0.049/0.667 \approx 0.024
\]
Half-life $\approx 29$ years. This is much closer to the empirical $\lambda\approx 0.02$ (half-life $\approx 35$ years). The inclusion of human capital explains the ``too-fast convergence'' puzzle of the basic Solow model (Mankiw, Romer, Weil, 1992).
\backbutton{q-v4}
\end{frame}

%% ----- SOLUTION VERY HARD 5 -----
\begin{frame}{Solution: Very Hard 5 (a)--(b)}
\hypertarget{s-v5}{}

\textbf{(a)} Without nPg, choose $c_t = M$ for all $t$ and finance via assets: $a_{t+1} = (1+r)a_t - M$. Starting from $a_0>0$, eventually $a_t$ becomes negative and grows as $a_t \sim -(1+r)^t M/r$. Debt explodes, but no constraint prevents it. Lifetime utility $\sum\beta^t u(M) = u(M)/(1-\beta)$, which grows without bound as $M\to\infty$. The problem has no finite solution---utility is unbounded above.

\medskip
\textbf{(b)} Iterate $c_t + a_{t+1} = (1+r)a_t$ forward:
\[
a_{T+1}/(1+r)^T = a_0(1+r) - \sum_{t=0}^{T}c_t/(1+r)^t
\]
Taking $T\to\infty$, nPg ($\lim a_{T+1}/(1+r)^T \geq 0$) implies:
\[
\sum_{t=0}^{\infty}\frac{c_t}{(1+r)^t} \leq a_0(1+r)
\]
\end{frame}

\begin{frame}{Solution: Very Hard 5 (c)--(d)}
\textbf{(c)} Suppose $\lim_{t\to\infty}\beta^t u'(c_t)a_{t+1} = L > 0$. This means assets retain positive shadow value at infinity. Consider a perturbation: consume slightly more at date $T$ by reducing $a_{T+1}$ by $\varepsilon$. Utility gain: $\beta^T u'(c_T)\varepsilon > 0$. Cost: lost continuation value $\approx \beta^{T+1}u'(c_{T+1})(1+r)\varepsilon$, which by Euler equals $\beta^T u'(c_T)\varepsilon$. But the TVC failure means the limit of these shadow values is positive---the household has ``leftover wealth'' at infinity it could consume. Hence the original path is suboptimal.

\medskip
\textbf{(d)} With $\beta(1+r)=1$: Euler gives $c_t = c_0$ for all $t$. From the budget: $c_0\sum_{t=0}^{\infty}\beta^t = a_0(1+r)$, so $c_0 = a_0(1-\beta)(1+r)$. Then $a_t = a_0$ for all $t$ (constant assets). TVC: $\lim_{t\to\infty}\beta^t c_0^{-\sigma} a_0 = a_0 c_0^{-\sigma}\lim\beta^t = 0$ since $\beta<1$. Holds if $a_0$ is finite.
\backbutton{q-v5}
\end{frame}

%% ----- SOLUTION VERY HARD 6 -----
\begin{frame}{Solution: Very Hard 6 (a)--(b)}
\hypertarget{s-v6}{}

\textbf{(a)} Linearize: $x_{t+1}-x^* \approx G'(x^*)(x_t-x^*)$, so $x_t-x^* = [G'(x^*)]^t(x_0-x^*)$.

If $G'(x^*)\in(0,1)$: $[G'(x^*)]^t > 0$ and shrinking $\implies$ deviation keeps same sign and vanishes. \textbf{Monotone convergence.}

If $G'(x^*)\in(-1,0)$: $[G'(x^*)]^t$ alternates sign but $|G'(x^*)|^t\to 0$ $\implies$ deviation oscillates with shrinking amplitude. \textbf{Oscillatory convergence.}

\medskip
\textbf{(b)} $G'(\bar{k}) = 1-\delta+s\alpha\bar{k}^{\alpha-1} = 1-\delta+\alpha\delta = 1-\delta(1-\alpha)$. For $\alpha\in(0,1)$ and $\delta\in(0,1)$: $\delta(1-\alpha)\in(0,1)$, so $G'(\bar{k})\in(0,1)$. Monotone convergence always holds. $G'(\bar{k})<0$ is impossible because both $1-\delta>0$ (for reasonable $\delta$) and $s\alpha\bar{k}^{\alpha-1}>0$.
\end{frame}

\begin{frame}{Solution: Very Hard 6 (c)--(d)}
\textbf{(c)} With $f(k)=k^\alpha e^{-\eta k}$: $G(k)=(1-\delta)k+sk^\alpha e^{-\eta k}$.
\[
G'(k) = 1-\delta+s(\alpha k^{\alpha-1}-\eta k^\alpha)e^{-\eta k}
\]
At steady state: $G'(\bar{k}) = 1-\delta+s\bar{k}^{\alpha-1}e^{-\eta\bar{k}}(\alpha-\eta\bar{k})$.

Using $s\bar{k}^{\alpha-1}e^{-\eta\bar{k}} = \delta$ (from $\delta\bar{k}=sf(\bar{k})$):
\[
G'(\bar{k}) = 1-\delta+\delta(\alpha-\eta\bar{k}) = 1-\delta(1-\alpha+\eta\bar{k})
\]
For $G'(\bar{k})<0$: need $\delta(1-\alpha+\eta\bar{k})>1$, i.e., $\eta\bar{k} > 1/\delta-1+\alpha$.

\medskip
\textbf{(d)} Chaos requires $|G'(\bar{k})|>1$, i.e., $G'(\bar{k})<-1$: $\delta(1-\alpha+\eta\bar{k})>2$. With $\delta=0.05$: need $\eta\bar{k}>39+\alpha$. This requires extreme parameter values---capital would need to be in a region of severely declining production. Empirically implausible for aggregate economies, confirming that chaotic macroeconomic dynamics from the Solow structure are a theoretical curiosity, not an empirical concern.
\backbutton{q-v6}
\end{frame}

\end{document}
